% ===================================================================
%  કોષ્ટક નં. ૫ — ઘર અને તેનું બાંધકામ.
%
%  Traduction, faite en 2026, en gujarati d'aujourd'hui, sur l'ido de
%  texto/io. Regles de la colonne : voir 10-kostak-01.tex. Une seule
%  numerotation. Le tableau est un DIALOGUE d'un bout a l'autre :
%  trente-six attributions de parole, le plus grand nombre du livret.
%  Trente-quatre paires annoncees par parigi.py pour ce seul tableau.
%
%  C'EST LE TABLEAU QUI AVAIT OBLIGE kolonoj.py A DISTINGUER LE RECIT
%  DU DIALOGUE, ET LE GUJARATI DONNE LA TROISIEME REPONSE DU RELEVE.
%  Le javanais a deux LEXIQUES paralleles, ngoko et krama, et un
%  enfant qui parle a son oncle doit changer de mots : huit regles
%  avaient du quitter la cle « mot » pour la cle « narracio ». Le
%  haoussa n'avait pas de niveaux du tout et marquait le rapport dans
%  le VOCATIF. Le gujarati a un troisieme systeme, ni lexical ni
%  vocatif mais PRONOMINAL : તું pour l'intime et l'inferieur, તમે
%  pour le respect, આપ pour le tres grand respect — et le pronom
%  gouverne le VERBE, qui s'accorde. Ioannes dit donc « તમે » a son
%  oncle et l'oncle repond « તું » au garcon, de bout en bout, sans
%  qu'un seul nom change. Les treize regles de cette colonne visent
%  des mots hindis et non un registre : aucune n'a a se deplacer, et
%  le fichier se controle comme les quinze autres.
%
%  LES METIERS DU BATIMENT SONT DES NOMS DE CASTE, ET C'EST LA
%  PREMIERE FOIS DU RELEVE. Le haoussa datait ses metiers par leur
%  morphologie — ma- pour ce qui precede le siecle, mai pour ce qui a
%  ete fabrique hier. Le gujarati les date autrement : les vieux
%  metiers portent le nom de la JATI qui les exercait de pere en
%  fils, et le nom du metier, le nom de la communaute et le nom de
%  famille sont le meme mot.
%
%    કડિયો (108)   le macon        સુથાર (115)   le charpentier
%    લુહાર (125)   le forgeron     સોની          l'orfevre
%    દરજી          le tailleur     મોચી          le cordonnier
%    કુંભાર        le potier       ધોબી          le blanchisseur
%    માળી (129)    le jardinier    વણકર          le tisserand
%
%  Ce sont aussi onze noms de famille gujaratis courants. Ce qui est
%  arrive apres ne peut plus entrer dans cette liste et prend un mot
%  libre ou un mot etranger : ઇજનેર l'ingenieur, કોન્ટ્રાક્ટર
%  l'entrepreneur (77), પ્લમ્બર, ઇલેક્ટ્રિશિયન (124).
%
%  ET « મિસ્ત્રી » EST PORTUGAIS, ce qui fait le quatrieme apres મેજ,
%  પાદરી et અનનાસ. Le mot vient de MESTRE, le maitre-ouvrier, et il
%  designe en gujarati l'homme de metier qualifie — celui qui n'est
%  pas d'une caste d'artisans mais qui sait faire. Un mot venu de Goa
%  pour nommer exactement ce que le systeme des jati ne prevoyait
%  pas.
%
%  LE CABINET D'AISANCES (20) EST « સંડાસ », ET LE RELEVE A
%  MAINTENANT TROIS FACONS DE LE NOMMER. Le haoussa disait « bayan
%  gida », derriere la maison — la PLACE. L'ido dit « latrino » — le
%  latin. Le gujarati dit સંડાસ, qui vient du persan et qui ne veut
%  rien dire d'autre : c'est un mot sans image, et c'est justement
%  pourquoi il a servi. Une langue qui doit nommer une chose dont on
%  ne parle pas prend volontiers un mot etranger, parce qu'un mot
%  etranger ne dit rien.
% ===================================================================

%%K t05-apar-1 apar
\VUsaut{6.0mm}
\VUtitre{15.0pt}{60}{કોષ્ટક નં. ૫}
\VUsaut{1.4mm}
\VUtitre{11.4pt}{0}{ઘર અને તેનું બાંધકામ.}
\VUsaut{2.2mm}

%%K t05-tit-1 sub
\VUpk{10.2pt}{40}{સારો ભેટો.}
\VUsaut{3.0mm}

%%K t05-01-1 p
\textsc{યોહાનેસ}. --- કાકા, મને જાણવાની ઘણી ઇચ્છા છે કે \VUgras{ઘર}
કેવી રીતે બંધાય છે.

%%K t05-01-2 p
\textsc{કાકા} (80). --- એ તો ઘણું સહેલું છે, મારા દોસ્ત, મારી સાથે આવ,
હું તને એ બતાવીશ જે શેઠ બેર્નાર્દુસ (79) બંધાવે છે અને જેમાં હું
રહીશ.

%%K t05-01-3 p
\textsc{યોહાનેસ}. --- તો આ ઘરનો \VUgras{નકશો} \textsuperscript{(76)}
કોણે દોર્યો ?

%%K t05-01-4 p
\textsc{કાકા}. --- \VUgras{સ્થપતિએ} \textsuperscript{(75)} ; તેમણે ઘણું
સ્પષ્ટ ચિત્ર રજૂ કર્યું જે મંજૂર થયું, અને \VUgras{કોન્ટ્રાક્ટરે}
\textsuperscript{(77)} કામ શરૂ કર્યું.

%%K t05-01-5 p
\textsc{યોહાનેસ}. --- તો કોન્ટ્રાક્ટર કોણ છે ?

%%K t05-01-6 p
\textsc{કાકા}. --- એ છે શેઠ બ્લાન્કો, તારા મિત્ર યાકોબુસના પિતા. તે
સ્થપતિ શેઠ રુફો સાથે મળીને કારીગરોને દોરે છે.

%%K t05-01-7 p
હં ! આ રહ્યા શેઠ બ્લાન્કો બરાબર વખતસર પોતાની \VUgras{માપપટ્ટી}
\textsuperscript{(78)} સાથે આવે છે, અને શેઠ રુફો પોતાના નકશા સાથે.

%%K t05-01-8 p
નમસ્કાર, સાહેબો. તમે કેમ છો ?

%%K t05-01-9 p
\textsc{શેઠ રુફો}. --- ઘણું સારું, આભાર. નમસ્કાર, યોહાનેસ, આ છોકરો
કેટલો મોટો થઈ ગયો !

%%K t05-01-10 p
\textsc{કાકા}. --- હા, ખરેખર, તે નાનો પુરુષ છે. તે ઘણો ગંભીર છે, તે જે
જુએ તે બધા વિશે તેને સમજાવવું પડે છે. આજે તે જાણવા માગે છે કે ઘર કેવી
રીતે બંધાય છે.

%%K t05-tit-2 sub
\VUpk{10.2pt}{40}{કારીગરો.}
\VUsaut{3.0mm}

%%K t05-01-11 p
\textsc{શેઠ બ્લાન્કો}. --- તો, મારી સાથે આવ, નાના દોસ્ત, હું તને
હમણાં જ સમજાવીશ, કેમ કે મને ભણવા માગતાં છોકરાં ઘણાં ગમે છે.

%%K t05-01-12 p
તું જુએ છે તે કારીગર જે હાથમાં \VUgras{પાવડો} \textsuperscript{(82)}
પકડે છે : તે \VUgras{ખોદનારો} \textsuperscript{(81)} છે. તે પાણીની
પાઇપ નાખવા માટે \VUgras{ખાઈ} \textsuperscript{(46)} ખોદે છે. તેણે ઘર
જ્યાં છે ત્યાંની જમીન પણ ખોદી છે, જેથી કડિયો ચાર \VUgras{ભીંતો} ચણી
શકે. આ ભીંતોનો જે ભાગ જમીનમાં છે તેને \VUgras{પાયો}
\textsuperscript{(2)} કહે છે. પાયાની વચ્ચેની જગ્યા \VUgras{ભોંયરું}
\textsuperscript{(3)} બનાવે છે. આ ભોંયરાને \VUgras{ભોંયરાની બારી}
\textsuperscript{(5)} નામની નાની બારી, \VUgras{બારણું}
\textsuperscript{(4)} અને દાદર છે.

%%K t05-01-13 p
\VUgras{કડિયાએ} \textsuperscript{(108)} \VUgras{બારણાં}
\textsuperscript{(8)} અને \VUgras{બારીઓ} \textsuperscript{(17)} માટે
ખાલી જગ્યા છોડીને ભીંતો ચણવાનું ચાલુ રાખ્યું.

%%K t05-01-14 p
\textsc{યોહાનેસ}. --- પણ એ કડિયો મને દેખાતો નથી !

%%K t05-01-15 p
\textsc{શેઠ બ્લાન્કો}. --- હમણાં તેણે ઘરમાં કામ પૂરું કર્યું છે. તે
\VUgras{ઘોડારની} \textsuperscript{(54)} પાસે છે, પોતાના
\VUgras{પાલખ} \textsuperscript{(110)} પર, અને \VUgras{ધોબીઘરની}
\textsuperscript{(56)} ભીંત ચણે છે.

%%K t05-01-16 p
તેની પાસે કરણી, ચૂનાનું તગારું અને \VUgras{ઓળંબો}
\textsuperscript{(109)} છે. પણ તું આ નામ યાદ રાખી શકીશ નહીં.

%%K t05-01-17 p
\textsc{યોહાનેસ}. --- જરૂર રાખીશ, કેમ કે હું હમણાં જ મારા કાકા પાસે
નાની કરણી અને નાનું તગારું માગીશ, અને હું જાતે દોરીથી ઓળંબો બનાવીશ.

%%K t05-tit-3 sub
\VUsaut{3.0mm}
\VUpk{10.2pt}{40}{બાંધકામનો સામાન.}
\VUsaut{3.0mm}

%%K t05-01-18 p
\textsc{કાકા}. --- વાહ ! ઘણું સારું. પણ મને કહે, યોહાનેસ, ઘર બાંધવા
માટે શું જોઈએ ?

%%K t05-01-19 p
\textsc{યોહાનેસ}. --- \VUgras{ઘડેલા પથ્થર} \textsuperscript{(59)} અને
\VUgras{ચૂનો} \textsuperscript{(65)} જોઈએ.

%%K t05-01-20 p
\textsc{કાકા}. --- બરાબર, પેલા મોટી ટોપીવાળા માણસને જો, પોતાના
\VUgras{ગાડા} \textsuperscript{(136)} પર. તે \VUgras{ગાડાવાળો}
\textsuperscript{(135)} છે. તે રોજ અહીં \VUgras{પથ્થરના પડ}
\textsuperscript{(60)} લાવે છે. તે પોતાનું ગાડું ચોકમાં ખાલી કરે છે,
અને \VUgras{પથ્થર ઘડનારા} \textsuperscript{(83)} પડ ઘડે છે અને પોતાના
\VUgras{પથ્થરકાપ કરવત} \textsuperscript{(85)} થી કાપે છે.
\VUgras{મજૂર} \textsuperscript{(146)} તેમને કડિયા પાસે લઈ જાય છે, જે
તેમને ગોઠવે છે.

%%K t05-01-21 p
એ દરમિયાન \VUgras{ચૂનો ગૂંદનારો} \textsuperscript{(87)} ચૂનો તૈયાર કરે
છે. તેને \VUgras{રેતી} \textsuperscript{(62)}, \VUgras{ચૂનો}
\textsuperscript{(63)} અને \VUgras{પાણી} \textsuperscript{(64)} જોઈએ,
જે \VUgras{કોથળામાં} \textsuperscript{(71)} તેની પાસે છે —
\VUgras{કડિયાનો મદદનીશ} \textsuperscript{(111)} તેને
\VUgras{હાથગાડી} \textsuperscript{(112)} પર રેતી લાવે છે, અને
\VUgras{શિખાઉ} \textsuperscript{(113)} છોકરો પંપ (58) પાસે છે. તે
\VUgras{ડોલ} \textsuperscript{(114)} ભરે છે. ચૂનો થોડી વારમાં તૈયાર થઈ
જશે. \VUgras{ચૂનો ઊંચકનારો} \textsuperscript{(107)} તેને લે છે અને
નિસરણી વડે પાલખ પર ઉપર લઈ જાય છે, જ્યાં કડિયો ઘરની આ બાજુ પૂરી કરે
છે.

%%K t05-01-22 p
અને હવે, જે કારીગર \VUgras{છાપરાનું} \textsuperscript{(31)} કામ કરે છે
તેને શું કહે છે ?

%%K t05-01-23 p
\textsc{યોહાનેસ}. --- તે \VUgras{છાપરાવાળો} \textsuperscript{(118)} છે.

%%K t05-tit-4 sub
\VUsaut{3.0mm}
\VUpk{10.2pt}{40}{ઘરનું છાપરું.}
\VUsaut{3.0mm}

%%K t05-01-24 p
\textsc{શેઠ બ્લાન્કો}. --- હા, પણ પહેલાં \VUgras{સુથાર}
\textsuperscript{(115)} આવે છે.

%%K t05-01-25 p
સુથાર \VUgras{છાપરાનું માળખું} \textsuperscript{(35)} બનાવે છે. તે
\VUgras{ભારવટિયા} \textsuperscript{(37)} જોડે છે \VUgras{ખીલા}
\textsuperscript{(117)} વડે. ઉપર જે કારીગરો છે, પોતાની પહોળી મખમલી
ચડ્ડી સાથે, તે સુથાર છે. એક નાનો ભારવટિયો પકડે છે, અને બીજો પોતાની
\VUgras{કુહાડીથી} \textsuperscript{(116)} ખીલો કાપે છે. જે
\VUgras{ધુમાડિયા} \textsuperscript{(41)} પાસે છે તે છાપરાવાળો છે. તે
(*) વળીઓ પર પાટિયાં જડે છે, અને પાટિયાં પર \VUgras{પથ્થરની તકતીઓ}
\textsuperscript{(66)}. ક્યારેક તે નળિયાં મૂકે છે.

%%K t05-noto-1 noto
\VUnotes{143.2mm}{%
(*) \VUgras{Balk-o} = D. \textit{Balken}, E. \textit{joist}, F.
\textit{solive}, I. \textit{travicello}, R. \textit{balka}, S. Port.
\textit{viga}.
અત્યાર સુધી નહીં સ્વીકારાયેલું ટેક્નિકલ મૂળ, પણ F. \textit{solive}
શબ્દનો અર્થ અનુવાદ કરવા માટે રજૂ કરવા જેવું, જે trabo નથી, F.
\textit{poutre}, કે trabeto પણ નથી. તે ચોરસ ઘડેલો લાકડાનો દાંડો છે જે
ભોંયતળિયું અને છત ટેકવે છે.%
}

%%K t05-01-26 p
\VUgras{ઘોડારનું છાપરું} \textsuperscript{(55)} નળિયાંનું છે,
\VUgras{ઘરોનું} \textsuperscript{(32)} પથ્થરની તકતીઓનું, અને
\VUgras{ઓસરીનું} \textsuperscript{(24)} જસતનું.

%%K t05-01-27 p
\textsc{યોહાનેસ}. --- શું છાપરાવાળો જસતનાં છાપરાંનું કામ પણ કરે છે ?

%%K t05-01-28 p
\textsc{શેઠ બ્લાન્કો}. --- ના, પણ \VUgras{જસતકામ કરનારો}
\textsuperscript{(121)} કે \VUgras{સીસાકામ કરનારો}, જે ઘરના છાપરા પર
\VUgras{છાપરાની બારી} \textsuperscript{(38)} મૂકે છે. તે
\VUgras{વરસાદની નળીઓ} \textsuperscript{(43)} અને \VUgras{છાપરાની
નીકો} \textsuperscript{(42)} પણ મૂકે છે. તે જસત અને પતરું સાંધે છે.
તેણે \VUgras{વીજરક્ષક} \textsuperscript{(40)} અને \VUgras{પવનસૂચક}
\textsuperscript{(39)} \VUgras{છાપરાના મોભ} \textsuperscript{(36)} પર
જડ્યાં છે.

%%K t05-01-29 p
\textsc{યોહાનેસ}. --- તો ઘર પૂરું થઈ ગયું ?

%%K t05-tit-5 sub
\VUsaut{3.0mm}
\VUpk{10.2pt}{40}{ઓજારો.}
\VUsaut{3.0mm}

%%K t05-01-30 p
\textsc{શેઠ બ્લાન્કો}. --- પૂરેપૂરું નહીં. \VUgras{પ્લાસ્ટર
કરનારો} \textsuperscript{(99)} \VUgras{ઈંટો} \textsuperscript{(61)} વડે
તે ભીંતો ચણે છે જે તેને જુદા જુદા ઓરડામાં વહેંચે છે. તે \VUgras{પ્લાસ્ટર}
\textsuperscript{(70)} લગાડે છે \VUgras{છત} \textsuperscript{(26)} અને
ભીંતો પર. હમણાં તે \VUgras{ઓસરીની}
\textsuperscript{(33)} છતનું કામ કરે છે. તેની પાસે, કડિયાની જેમ,
\VUgras{કરણી} \textsuperscript{(100)} અને \VUgras{તગારું}
\textsuperscript{(101)} છે.

%%K t05-01-31 p
પછી કારીગર બારણાં, બારીઓ, \VUgras{લાકડાનું ભોંયતળિયું}
\textsuperscript{(16)} અને \VUgras{બારીનાં કમાડ} \textsuperscript{(19)}
બનાવે છે.

%%K t05-01-32 p
હમણાં તે ચોકમાં પોતાના \VUgras{કામના પાટિયા} \textsuperscript{(90)} પર
કામ કરે છે. તેની પાસે \VUgras{કરવત} \textsuperscript{(94)} છે લાકડું
(69) કાપવા માટે, \VUgras{રંધો} \textsuperscript{(91)},
\VUgras{ભીંસ} \textsuperscript{(92)}, \VUgras{ટાંકણું}
\textsuperscript{(94 bis)}, \VUgras{પરિકર} \textsuperscript{(95)}, અને
લાકડાની \VUgras{હથોડી} \textsuperscript{(95 bis)}.

%%K t05-01-33 p
\textsc{યોહાનેસ}. --- વાહ ! પણ સુથારીકામ કડિયાકામ કરતાં વધારે મજાનું
છે. કાકા, તમે મારા માટે કામનું પાટિયું, કરવત અને રંધો ખરીદશો, કેમ કે
હું થોડી વારમાં મેદોર માટે ઝૂંપડી બનાવીશ.

%%K t05-01-34 p
\textsc{કાકા}. --- હા, નાના.

%%K t05-01-35 p
\textsc{યોહાનેસ}. --- હું કેટલો ખુશ છું ! અને પ્રવેશદ્વાર પાસે જે
ઊભો છે તે કોણ છે ?

%%K t05-tit-6 sub
\VUsaut{3.0mm}
\VUpk{10.2pt}{40}{રંગકામ.}
\VUsaut{3.0mm}

%%K t05-01-36 p
\textsc{શેઠ બ્લાન્કો}. --- તે \VUgras{રંગારો} \textsuperscript{(102)}
છે. તે પોતાની નિસરણી પર ઊભો છે \VUgras{રંગના}
\textsuperscript{(103)} ડબા અને પોતાના \VUgras{પીંછી}
\textsuperscript{(104)} સાથે. તે પ્રવેશદ્વારનું લાકડું વધારે વખત ટકે
તે માટે રંગે છે.

%%K t05-01-37 p
તે ઓરડાઓમાં કાગળ પણ ચોંટાડે છે.

%%K t05-01-38 p
\VUgras{સજાવટ કરનારો} \textsuperscript{(107)}, જે \VUgras{નિસરણી}
\textsuperscript{(106)} પર છે, \VUgras{ઝરૂખા} \textsuperscript{(28)} પર,
\VUgras{પડદો} \textsuperscript{(29)} મૂકે છે.

%%K t05-01-39 p
મોટી બારી પાસે જે કારીગર ઊભો છે તે \VUgras{કાચ બેસાડનારો}
\textsuperscript{(96)} છે. તે \VUgras{કાચની તકતીઓ}
\textsuperscript{(18)} બેસાડે છે \VUgras{ચીકણા લેપ} \textsuperscript{(73)}
વડે. બીજો કારીગર, જે ભોંયરાના બારણા પાસે ઘૂંટણિયે પડ્યો છે, તે
\VUgras{તાળાં બનાવનારો} \textsuperscript{(97)} છે. તે \VUgras{તાળું}
\textsuperscript{(6)} બેસાડે છે. તેની \VUgras{કાનસ}
\textsuperscript{(98)} તેની પાસે છે.

%%K t05-01-40 p
\textsc{યોહાનેસ}. --- જે કારીગર લોખંડ પર પોતાની \VUgras{હથોડી}
\textsuperscript{(84)} મારે છે તે પણ તાળાં બનાવનારો છે ?

%%K t05-01-41 p
\textsc{શેઠ બ્લાન્કો}. --- વધારે બરાબર કહીએ તો, તે લુહાર છે (125). તે
\VUgras{લોખંડનો સામાન} \textsuperscript{(67)} ઘડે છે
\VUgras{વીજકારીગર} \textsuperscript{(124)} માટે, જે \VUgras{ગાડીખાનાના}
\textsuperscript{(51)} છાપરા પર છે. તેની પાસે \VUgras{ભઠ્ઠી}
\textsuperscript{(126)}, \VUgras{એરણ} \textsuperscript{(127)} અને
\VUgras{સાણસી} \textsuperscript{(128)} છે.

%%K t05-01-42 p
વીજકારીગર ટેલિફોનનો \VUgras{તાંબાનો} \textsuperscript{(44)}
\VUgras{તાર} બેસાડે છે.

%%K t05-tit-7 sub
\VUpk{10.2pt}{40}{બગીચાનું કામ.}
\VUsaut{3.0mm}

%%K t05-01-43 p
\textsc{કાકા}. --- \VUgras{ચોકના} \textsuperscript{(45)} છેડે,
\VUgras{જાળી} \textsuperscript{(47)} અને \VUgras{મોટા દરવાજા}
\textsuperscript{(48)} પાસે તું \VUgras{માળીને} \textsuperscript{(129)}
જુએ છે. તે \VUgras{નાનો બગીચો} \textsuperscript{(49)} તૈયાર કરે છે.
તેણે \VUgras{કોદાળીથી} \textsuperscript{(131)} જમીન ખોદી, તે બી વાવે
છે અને \VUgras{દાંતાળીથી} \textsuperscript{(130)} સરખું કરે છે. પછી,
પોતાના \VUgras{ઝારાથી} \textsuperscript{(132)}, તે \VUgras{ક્યારાને}
\textsuperscript{(150)} પાણી પાશે અને છોડ ઊગશે.

%%K t05-01-44 p
\VUgras{ઘોડાનો સાઈસ} \textsuperscript{(133)}, જેણે હમણાં ઘોડાની
સંભાળ લીધી અને તેને ખાવાનું આપ્યું, તે \VUgras{ગાડી}
\textsuperscript{(52)} સાફ કરે છે સ્પોન્જ, \VUgras{હાથપંપ}
\textsuperscript{(134)} અને પાણીની ડોલથી. પછી તે તેને ગાડીખાનામાં પાછી
મૂકશે અને જાડા \VUgras{આગળિયાથી} \textsuperscript{(53)} બારણું બંધ
કરશે.

%%K t05-01-45 p
લે, તું ખુશ છે, હું માનું છું, તને પૂરતી સમજ મળી.

%%K t05-01-46 p
\textsc{યોહાનેસ}. --- હા, કાકા, પણ મને ઘરની અંદરનો ભાગ જોવાની ઇચ્છા
છે.

%%K t05-01-47 p
\textsc{કાકા}. --- એ કાલે થશે. શેઠ રુફો તને નકશો સમજાવશે અને દરેક
ઓરડાનું નામ બતાવશે.

%%K t05-tit-8 sub
\VUsaut{3.0mm}
\VUpk{10.2pt}{40}{ઘરની અંદરની ગોઠવણ.}
\VUsaut{2.4mm}

%%K t05-apar-2 apar
{\centering\textit{(નકશો જુઓ.)}\par}
\VUsaut{2.4mm}

%%K t05-01-48 p
\textsc{સ્થપતિ}. --- આવ, યોહાનેસ, હું તને હમણાં જ ઘરના જુદા જુદા ભાગ
બતાવીશ.

%%K t05-01-49 p
આપણે \VUgras{ભોંયતળિયામાં} \textsuperscript{(7)} જઈએ. આ રહ્યું
\VUgras{ઘંટડીનું} \textsuperscript{(11)} બટન. આપણે \VUgras{પગથિયાંનાં}
\textsuperscript{(14)} ત્રણ \VUgras{પગથિયાં} ચઢીએ છીએ, જે
\VUgras{ઓટલાનાં} \textsuperscript{(9)} છે, આપણે \VUgras{ઉંબરો}
\textsuperscript{(10)} ઓળંગીએ છીએ, જે \VUgras{પ્રવેશદ્વારનો}
\textsuperscript{(8)} છે, અને આ રહ્યા આપણે \VUgras{દાદરના}
\textsuperscript{(13)} પગથિયે.

%%K t05-01-50 p
\textsc{યોહાનેસ}. --- આ ગલી છે.

%%K t05-01-51 p
\textsc{સ્થપતિ}. --- ના, આ \VUgras{પ્રવેશખંડ} \textsuperscript{(12)}
છે. \VUgras{ગલી} \textsuperscript{(a)} છેડે છે. જમણી બાજુએ, આ રહ્યો
\VUgras{બેઠકખંડ} \textsuperscript{(b)} અને \VUgras{ભોજનખંડ}
\textsuperscript{(c)}, ભોજનખંડની સામે \VUgras{રસોડું}
\textsuperscript{(d)} અને \VUgras{નોકરખંડ} \textsuperscript{(e)} છે, પછી
આગળ \VUgras{લેખનખંડ} \textsuperscript{(f)}. \VUgras{ઓસરી}
\textsuperscript{(23)} પોતાના \VUgras{કાચના બારણા}
\textsuperscript{(25)} સાથે બેઠકખંડ પાસે છે.

%%K t05-01-52 p
આપણે હમણાં દાદર ચઢીશું. \VUgras{કઠેડો} \textsuperscript{(15)} પકડ અને
પગથિયાં ગણ. તે ચૌદ છે. આ રહી \VUgras{પરસાળ} \textsuperscript{(m)},
આપણે \VUgras{પહેલા માળે} \textsuperscript{(27)} છીએ.

%%K t05-tit-9 sub
\VUsaut{3.0mm}
\VUpk{10.2pt}{40}{પહેલો માળ.}
\VUsaut{3.0mm}

%%K t05-01-53 p
દાદરની સામે \VUgras{સંડાસ} \textsuperscript{(20)} અને \VUgras{શણગારખંડ}
\textsuperscript{(j)} છે. જમણી બાજુએ સુંદર \VUgras{શયનખંડ}
\textsuperscript{(g)} છે, જેને ઘરના \VUgras{આગલા ભાગ}
\textsuperscript{(1)} પર બે બારી છે. ડાબી બાજુએ \VUgras{સ્નાનખંડ}
\textsuperscript{(1)} અને \VUgras{મહેમાનખંડ} \textsuperscript{(i)} છે,
જેને મોટો \VUgras{ગોખલો} \textsuperscript{(h)} છે : શયનખંડ પાસે
\VUgras{બાળકોનો ખંડ} \textsuperscript{(k)} છે. તે વિશાળ \VUgras{અગાસી}
\textsuperscript{(21)} પાસે છે. આ અગાસી નીચે બીજું સંડાસ (20) છે, અને
ભીંતે, આ રહ્યો \VUgras{ઘંટ} \textsuperscript{(22)}, જે રાતના ભોજન માટે
વાગે છે.

%%K t05-01-54 p
\textsc{યોહાનેસ}. --- આપણે \VUgras{બીજા માળે} જઈએ.

%%K t05-01-55 p
\textsc{સ્થપતિ}. --- બીજો માળ છે જ નહીં. છાપરા નીચે માત્ર
\VUgras{માળિયાના ઓરડા} \textsuperscript{(33)} અને અનાજનો કોઠાર (34) છે.
આપણે તપાસ પૂરી કરી, આપણે નીચે ઊતરી શકીએ.

%%K t05-01-56 p
\textsc{યોહાનેસ}. --- આભાર, સાહેબ, આ ઘણું રસપ્રદ છે. હું હમણાં જ મારા
પિતાને મેં જે જોયું તે બધું કહીશ.

\VUsaut{4.0mm}
\VUfilet{22mm}
